\chapter{Смешанная полярно-передаточная кривизна связывающего бимодуля}
\label{ch:v7-polar-transfer-cross-curvature-origin-gate}

\section{Замена вопроса о весе вопросом о кривизне}

Предыдущий гейт запретил сохранять сырую норму
$\frac12\|Z\|_{\rm HS}^2$ после диагонального quotient: метрика двух
исходных углов индуцирует $\|Z\|_{\rm HS}^2$ и возвращает неверную
сигнатуру. Однако связывающая алгебра допускает другой канонический объект.
Она не обязана считать среднее двух кривизн; она несёт кривизну
эквивалентного бимодуля, то есть их ковариантную разность.

Пусть
\begin{equation}
 A_0:\mathbb C^{11}\longrightarrow\mathbb C^{10},
 \qquad A_0=U|A_0|
 \label{eq:v7-polar-linking-reference}
\end{equation}
--- полярное разложение фоновой инцидентности. Тогда
\begin{equation}
 UU^*=I_{10},\qquad U^*U=P,\qquad \rank P=10,
 \qquad \rank(I_{11}-P)=1.
 \label{eq:v7-polar-linking-support}
\end{equation}
Для переменного прямоугольного поля $A$ определим две относительные
Gram-кривизны
\begin{equation}
 C_s(A)=A^*A-A_0^*A_0,
 \qquad
 C_t(A)=AA^*-A_0A_0^*.
 \label{eq:v7-polar-linking-end-curvatures}
\end{equation}

\section{Канонический знак разности}

Полярная часть $U$ является элементом прямоугольного связывающего
бимодуля. Индуцированная на нём кривизна равна
\begin{equation}
 \mathcal R_U(A)=C_t(A)U-UC_s(A).
 \label{eq:v7-polar-linking-relative-curvature}
\end{equation}
Знак минус не является параметром: он следует из правого правила Лейбница
для двусторонней связности
\begin{equation}
 \nabla\xi=d\xi+A_t\xi-\xi A_s.
 \label{eq:v7-polar-linking-bimodule-connection}
\end{equation}
Это тот же механизм относительной моритовой кривизны, который был независимо
зафиксирован в Томе~V.

Самосопряжённое связывающее завершение имеет вид
\begin{equation}
 \mathbb F_U=
 \begin{pmatrix}
  0&\mathcal R_U\\
  \mathcal R_U^*&0
 \end{pmatrix}.
 \label{eq:v7-polar-linking-selfadjoint-completion}
\end{equation}
Поэтому
\begin{equation}
 \frac14\|\mathbb F_U\|_{\rm HS}^2
 =\frac12\|\mathcal R_U\|_{\rm HS}^2.
 \label{eq:v7-polar-linking-half-trace-identity}
\end{equation}
Здесь множитель $1/2$ удаляет две ориентированные записи одного
внедиагонального блока. Он не является относительным весом между рёберным и
физическим секторами.

\section{Почему опасный нулевой гессиан исчезает}

Полярное разложение даёт точное переплетение фоновых углов:
\begin{equation}
 (A_0A_0^*)U=U(A_0^*A_0).
 \label{eq:v7-polar-linking-background-intertwining}
\end{equation}
В точке $A=0$ обе кривизны имеют только постоянную фоновую часть, и она
сокращается в~\eqref{eq:v7-polar-linking-relative-curvature}:
\begin{equation}
 \mathcal R_U(0)=0.
 \label{eq:v7-polar-linking-zero-background}
\end{equation}
Поскольку $C_s(A)-C_s(0)$ и $C_t(A)-C_t(0)$ квадратичны по $A$, имеем
\begin{equation}
 \mathcal R_U(A)=O(A^2),
 \qquad
 \Hess_0\frac12\|\mathcal R_U(A)\|_{\rm HS}^2=0.
 \label{eq:v7-polar-linking-zero-hessian}
\end{equation}
Следовательно, общий локальный функционал
\begin{equation}
 \mathcal S_{\rm link}(A)
 =\mathcal S_E(A)+\frac12\|\mathcal R_U(A)\|_{\rm HS}^2
 \label{eq:v7-polar-linking-action}
\end{equation}
наследует в нуле ровно рёберный запуск:
\begin{equation}
 \operatorname{SpecSign}\Hess_0\mathcal S_{\rm link}
 =(7,0,20),
 \qquad
 \lambda_{\rm heavy}^{\min}=\frac{18}{5}.
 \label{eq:v7-polar-linking-origin-signature}
\end{equation}
Семь выбранных корней неустойчивы, а двадцать конкурирующих направлений
имеют положительную щель без настройки прежнего коэффициента $\beta$.

\section{Полный гессиан выбранного вакуума}

В точке $A=A_0$ обе относительные Gram-кривизны равны нулю. Линеаризация
связывающей кривизны имеет вид
\begin{align}
 \delta\mathcal R_U
 &=(A_0\delta A^*+\delta A A_0^*)U\nonumber\\
 &\quad-U(A_0^*\delta A+\delta A^*A_0).
 \label{eq:v7-polar-linking-linearization}
\end{align}
Её квадрат является положительной полуопределённой добавкой ранга $22$ с
пятью нулевыми направлениями. Рёберный Hodge-блок заполняет оставшиеся
нули. Полный 27-мерный результат равен
\begin{equation}
 \operatorname{SpecSign}\Hess_{A_0}\mathcal S_{\rm link}
 =(0,0,27),
 \qquad
 \lambda_{\min}=3.9368554658\ldots.
 \label{eq:v7-polar-linking-vacuum-signature}
\end{equation}

Индексный дефект не вставляется вручную. Полярная карта удовлетворяет
$U(I-P)=0$, а смешанные вариации между опорой и дефектной линией входят в
линеаризацию~\eqref{eq:v7-polar-linking-linearization}. Именно поэтому
полная прямоугольная кривизна имеет ранг $22$, тогда как её сжатие только к
согласованному углу теряло одно дополнительное направление.

\section{Контроль ложных альтернатив}

Знак суммы
\begin{equation}
 \mathcal R_U^{(+)}=C_tU+UC_s
 \label{eq:v7-polar-linking-plus-channel}
\end{equation}
не является кривизной правого модуля. Он сохраняет общий фоновый угол и при
той же связывающей нормировке даёт сигнатуру нуля
\begin{equation}
 (27,0,0),
 \label{eq:v7-polar-linking-plus-fail}
\end{equation}
то есть делает неустойчивыми уже все направления.

Произвольная фаза между двумя концами также не выводится. Смешанный член
$\langle Z,D\rangle$ нечётен при обмене source/target и сокращается в
Real-симметричном завершении. Поэтому положительный результат не является
скрытой настройкой фазы или коэффициента.

\begin{theorem}[Локальный полярно-связывающий родитель]
Каноническая относительная кривизна
$\mathcal R_U=C_tU-UC_s$ существует на уже выведенном полярном
связывающем бимодуле $10\times11$. Её нулевой гессиан тождественно равен
нулю, поскольку фоновая инцидентность переплетает два Gram-угла. В сумме с
рёберной Hodge-нормой она даёт сигнатуру $(7,0,20)$ в нуле и строго
положительный 27-мерный гессиан целевого вакуума. Ручной вес $\beta=1/2$
для селективного запуска больше не требуется.
\end{theorem}

\begin{remark}
Гейт является локальным допуском архитектуры, а не полным физическим
замыканием. Следующий тест обязан получить одновременно рёберную Hodge-
кривизну и $\mathcal R_U$ как блоки одной Real-суперсвязности с одним
следом. До этого количественные отношения масс и нормировка кинетики не
считаются выведенными.
\end{remark}

\begin{nogocriterion}[Запрет возврата к суммарному Gram-каналу]
Нельзя после положительного относительного прохода добавлять прежнюю сумму
$\|C_s\|^2+\|C_t\|^2$ или знак плюс и затем компенсировать её свободным
весом. Это восстанавливает уже закрытый гессиан $(21,0,6)$ или усиливает
провал до $(27,0,0)$.
\end{nogocriterion}

\begin{protocol}
\begin{enumerate}
 \item полярная коизометрия $11\to10$: \textbf{сохранена};
 \item индексный дефект ранга один: \textbf{сохранён};
 \item знак относительной кривизны: \textbf{выведен правилом Лейбница};
 \item ручной относительный вес в нулевом гессиане: \textbf{не требуется};
 \item сигнатура нуля: \textbf{$(7,0,20)$};
 \item тяжёлая щель нуля: \textbf{$18/5$};
 \item ранг связывающего вакуумного гессиана: \textbf{$22$};
 \item полный вакуум: \textbf{$(0,0,27)$};
 \item знак суммы: \textbf{провал $(27,0,0)$};
 \item количественная нормировка масс: \textbf{открыта};
 \item следующий гейт: \textbf{единая Real-связывающая суперсвязность}.
\end{enumerate}
\end{protocol}

Машинный сертификат сохранён в
\path{s2t_v7_polar_transfer_cross_curvature_origin_gate_results.json}.

\begin{thebibliography}{9}
\bibitem{v7-polar-linking-quillen}
D.~Quillen,
\emph{Superconnections and the Chern Character},
Topology 24 (1985), 89--95.
\bibitem{v7-polar-linking-kodaka-teruya}
K.~Kodaka, T.~Teruya,
\emph{The Strong Morita Equivalence for Inclusions of $C^*$-Algebras and
Conditional Expectations for Equivalence Bimodules},
arXiv:1609.08263.
\bibitem{v7-polar-linking-project}
Том~V, глава о связывающем родителе Мориты и двусторонней кривизне.
\end{thebibliography}